\documentclass[../../main.tex]{subfiles}

\begin{document}

\section{Obbiettivi futuri} \label{sec:obbiettivi-futuri}
  %    - capacità di creare ricostruzioni verosimili prima di aver finito di
  %    implementare i vari algoritmi di ricostruzione veri e propri
  %    - capacità di creare oggetti in output completi anche in assenza di una
  %    ricostruzione completa
  Al momento della stesura di questa tesi fake-reco non implementa le seguenti
  funzionalità ma sono previste dalla \textit{roadmap}:

  \paragraph{\textit{Smearing} della ricostruzione}
    La prima miglioria da apportare allo stato attuale delle cose è aggiungere
    uno \textit{smearing} ai dati inseriti nei campi della ricostruzione, così
    da rendere più realistici i dati sui quali testare i \textit{workflow} di
    analisi; con questa modifica il nome passa da fake-reco a fast-reco.

    \begin{figure}[ht]
      \centering
      \includegraphics[width=.95\textwidth]{fake-reco/fake-recov2}
      \caption{
        Schema di un \textit{workflow} che utilizza la fast-reco divisa in
        moduli specifici per i singoli detector.
      }
      \label{fig:fake-recov2}
    \end{figure}

    \begin{figure}[ht]
      \centering
      \includegraphics[width=.95\textwidth]{fake-reco/not-that-fake-reco}
      \caption{
        Schema di un \textit{workflow} di una possibile fase dello sviluppo che
        affianca alcuni degli effettivi moduli di ricostruzione a quelli della
        fast-reco per i detector che ancora non ne hanno uno.
      }
      \label{fig:not-that-fake-reco}
    \end{figure}

  \paragraph{fake-reco parziale}
    L'ultimo step prima di abbandonare fast-reco in favore delle ricostruzioni
    vere e proprie, schematizzato in figura~\ref{fig:fake-recov2}, consiste
    nel dividerlo in quattro processi più piccoli: la traduzione della verità
    Monte Carlo (non indicato come modulo nello schema) e tre moduli di
    fast-reco ma specifici per ogni componente di SAND\@.
    In questo modo i veri algoritmi di ricostruzione, che non saranno
    portati a termine tutti contemporaneamente, potranno essere inseriti uno
    alla volta, continuando ad avere in output degli oggetti completi; una
    possibile fase intermedia di questo processo è schematizzata in
    figura~\ref{fig:not-that-fake-reco}.

\end{document}